% Changed: restructure for a single evidence-bound service-object argument.
% Protected: measured values, protocol distinctions and failure records.
% Meaning-Check: PRESERVED
% Risk-Flags: none
\begin{abstract}
A retrieval service can repeat thousands of document tokens while changing a question placed before them. Native prefix caching then misses most of the repeated context. Reordering exposes a longer prefix, but can change model output and propagate unsupported state across tenants. We present ContractKV, which treats a compiled KV prefix as a service object admitted by scoped behavioral evidence and made visible through a transaction. Contract-Aware Layout Compilation preserves registered dependencies; Scoped Behavioral Certification checks a model-generated token path and termination; Prefix-Selective Transactional Promotion commits only the exact, block-aligned stable prefix; and Service-Template Certificates bind authorization and enable qualification costs to be amortized. We distinguish deployment-qualified templates from a stricter query-qualified protocol whose first raw business response supplies the witness. On six source-disjoint HotpotQA and RepoBench templates, three fresh-process repetitions yield full-cost throughput gains of 1.619$\times$, 1.109$\times$, and 1.070$\times$ for Qwen3-8B, Mistral-7B, and Llama-3.1-8B. All 1,296 admitted compiled future pairs retain their raw output token IDs. These zero-slack measurements include successful and failed qualification and waiting; p99 does not uniformly improve. A separate Qwen cohort yields 1.803$\times$. Live rejection, template-amortization, and cache-lifecycle experiments connect the gains to controlled publication and physical reuse. The results show where behavior-scoped compilation can turn prompt repetition into a measurable service benefit.
\end{abstract}
